\documentclass[12pt,a4paper]{article}
\usepackage[utf8]{vietnam}
\usepackage{amsmath,amssymb}
\usepackage{tabularx}
\usepackage{tikz}
\usepackage{xcolor}
\usepackage[top=0.8cm, bottom=0.8cm, left=1cm, right=1cm]{geometry}

\setlength{\parskip}{3pt}

\begin{document}
\pagestyle{plain}

% #############################################################################
% ========================= HƯỚNG DẪN CHẤM MÃ ĐỀ 901 ==========================
% #############################################################################
\noindent
\begin{tabularx}{\linewidth}{|X|c|}
\hline
\textbf{ĐÁP ÁN \& HƯỚNG DẪN CHẤM CHI TIẾT -- TOÁN 9} & \textbf{Mã đề: 901} \\
\hline
\end{tabularx}

\vspace{3mm}

\noindent\textbf{\large Bài 1 (4,0 điểm). Giải các hệ phương trình sau bằng phương pháp cộng đại số:}

\vspace{1mm}
\noindent\textbf{a) (2,0 điểm)} Giải hệ phương trình: $\begin{cases} 3x - 2y = 7 \quad &(1) \\ 2x + 2y = 8 \quad &(2) \end{cases}$
\begin{itemize}
	\item Nhận thấy hệ số của $y$ trong hai phương trình đối nhau ($-2$ và $2$).\\
	Cộng từng vế hai phương trình (1) và (2) ta được:
	$$(3x - 2y) + (2x + 2y) = 7 + 8 \iff 5x = 15 \iff x = 3. \quad \text{\itshape (1,0 điểm)}$$
	\item Thay $x = 3$ vào phương trình (2) ta được:
	$$2 \cdot 3 + 2y = 8 \iff 6 + 2y = 8 \iff 2y = 2 \iff y = 1. \quad \text{\itshape (0,75 điểm)}$$
	\item Kết luận: Vậy hệ phương trình đã cho có nghiệm duy nhất là $(x; y) = (3; 1)$. \quad \text{\itshape (0,25 điểm)}
\end{itemize}

\vspace{1.5mm}
\noindent\textbf{b) (2,0 điểm)} Giải hệ phương trình: $\begin{cases} 3x + 5y = 1 \quad &(1) \\ 2x - 3y = 7 \quad &(2) \end{cases}$
\begin{itemize}
	\item Nhân hai vế phương trình (1) với $3$ và nhân hai vế phương trình (2) với $5$, ta được hệ tương đương:
	$$\begin{cases} 9x + 15y = 3 \quad &(3) \\ 10x - 15y = 35 \quad &(4) \end{cases} \quad \text{\itshape (0,5 điểm)}$$
	\item Cộng từng vế của phương trình (3) và phương trình (4) ta được:
	$$(9x + 15y) + (10x - 15y) = 3 + 35 \iff 19x = 38 \iff x = 2. \quad \text{\itshape (0,75 điểm)}$$
	\item Thay $x = 2$ vào phương trình (1):
	$$3 \cdot 2 + 5y = 1 \iff 6 + 5y = 1 \iff 5y = -5 \iff y = -1. \quad \text{\itshape (0,5 điểm)}$$
	\item Kết luận: Vậy hệ phương trình đã cho có nghiệm duy nhất là $(x; y) = (2; -1)$. \quad \text{\itshape (0,25 điểm)}
\end{itemize}

\vspace{2mm}
\noindent\textbf{\large Bài 2 (3,0 điểm). Rút gọn biểu thức chứa căn thức:}

Cho biểu thức $P = \left( \dfrac{\sqrt{x}}{\sqrt{x}-1} - \dfrac{1}{x-\sqrt{x}} \right) : \left( \dfrac{1}{\sqrt{x}+1} + \dfrac{2}{x-1} \right)$ với $x > 0$ và $x \ne 1$.
\begin{itemize}
	\item Xét biểu thức trong ngoặc thứ nhất:
	$$A = \dfrac{\sqrt{x}}{\sqrt{x}-1} - \dfrac{1}{\sqrt{x}(\sqrt{x}-1)} = \dfrac{x - 1}{\sqrt{x}(\sqrt{x}-1)} = \dfrac{(\sqrt{x}-1)(\sqrt{x}+1)}{\sqrt{x}(\sqrt{x}-1)} = \dfrac{\sqrt{x}+1}{\sqrt{x}}. \quad \text{\itshape (1,25 điểm)}$$
	\item Xét biểu thức trong ngoặc thứ hai:
	$$B = \dfrac{1}{\sqrt{x}+1} + \dfrac{2}{(\sqrt{x}-1)(\sqrt{x}+1)} = \dfrac{\sqrt{x}-1+2}{(\sqrt{x}-1)(\sqrt{x}+1)} = \dfrac{\sqrt{x}+1}{(\sqrt{x}-1)(\sqrt{x}+1)} = \dfrac{1}{\sqrt{x}-1}. \quad \text{\itshape (1,0 điểm)}$$
	\item Thực hiện phép chia:
	$$P = A : B = \dfrac{\sqrt{x}+1}{\sqrt{x}} : \dfrac{1}{\sqrt{x}-1} = \dfrac{\sqrt{x}+1}{\sqrt{x}} \cdot (\sqrt{x}-1) = \dfrac{x-1}{\sqrt{x}}.$$
	Vậy $P = \dfrac{x-1}{\sqrt{x}}$. \quad \text{\itshape (0,75 điểm)}
\end{itemize}

\vspace{2mm}
\noindent\textbf{\large Bài 3 (3,0 điểm). Giải bài toán bằng cách lập hệ phương trình:}
\begin{itemize}
	\item Gọi số chi tiết máy tổ I phải làm theo kế hoạch là $x$ (chi tiết máy).\\
	Gọi số chi tiết máy tổ II phải làm theo kế hoạch là $y$ (chi tiết máy).\\
	Điều kiện: $x, y \in \mathbb{N}^*$ và $x, y < 900$. \quad \text{\itshape (0,5 điểm)}
	\item Vì theo kế hoạch hai tổ làm $900$ chi tiết máy nên ta có:
	$$x + y = 900 \quad (1) \quad \text{\itshape (0,5 điểm)}$$
	\item Thực tế tổ I vượt $15\%$ làm thêm $0{,}15x$; tổ II vượt $10\%$ làm thêm $0{,}1y$.\\
	Cả hai tổ làm được $1010$ chi tiết máy (vượt mức $1010 - 900 = 110$ chi tiết máy):
	$$0{,}15x + 0{,}1y = 110 \iff 3x + 2y = 2200 \quad (2) \quad \text{\itshape (0,75 điểm)}$$
	\item Giải hệ $\begin{cases} x + y = 900 \\ 3x + 2y = 2200 \end{cases} \iff \begin{cases} 2x + 2y = 1800 \\ 3x + 2y = 2200 \end{cases} \iff \begin{cases} x = 400 \\ y = 500 \end{cases}$ (thỏa mãn điều kiện). \quad \text{\itshape (0,75 điểm)}
	\item Kết luận: Theo kế hoạch tổ I làm $400$ chi tiết máy, tổ II làm $500$ chi tiết máy. \quad \text{\itshape (0,5 điểm)}
\end{itemize}

\newpage
% #############################################################################
% ========================= HƯỚNG DẪN CHẤM MÃ ĐỀ 902 ==========================
% #############################################################################
\noindent
\begin{tabularx}{\linewidth}{|X|c|}
\hline
\textbf{ĐÁP ÁN \& HƯỚNG DẪN CHẤM CHI TIẾT -- TOÁN 9} & \textbf{Mã đề: 902} \\
\hline
\end{tabularx}

\vspace{3mm}

\noindent\textbf{\large Bài 1 (4,0 điểm). Giải các hệ phương trình sau bằng phương pháp cộng đại số:}

\vspace{1mm}
\noindent\textbf{a) (2,0 điểm)} Giải hệ phương trình: $\begin{cases} 2x - 3y = 5 \quad &(1) \\ 3x + 3y = 15 \quad &(2) \end{cases}$
\begin{itemize}
	\item Nhận thấy hệ số của $y$ trong hai phương trình đối nhau ($-3$ và $3$).\\
	Cộng từng vế hai phương trình (1) và (2) ta được:
	$$(2x - 3y) + (3x + 3y) = 5 + 15 \iff 5x = 20 \iff x = 4. \quad \text{\itshape (1,0 điểm)}$$
	\item Thay $x = 4$ vào phương trình (2) ta được:
	$$3 \cdot 4 + 3y = 15 \iff 12 + 3y = 15 \iff 3y = 3 \iff y = 1. \quad \text{\itshape (0,75 điểm)}$$
	\item Kết luận: Vậy hệ phương trình đã cho có nghiệm duy nhất là $(x; y) = (4; 1)$. \quad \text{\itshape (0,25 điểm)}
\end{itemize}

\vspace{1.5mm}
\noindent\textbf{b) (2,0 điểm)} Giải hệ phương trình: $\begin{cases} 2x + 3y = 12 \quad &(1) \\ 3x - 2y = 5 \quad &(2) \end{cases}$
\begin{itemize}
	\item Nhân hai vế phương trình (1) với $2$ và nhân hai vế phương trình (2) với $3$, ta được hệ tương đương:
	$$\begin{cases} 4x + 6y = 24 \quad &(3) \\ 9x - 6y = 15 \quad &(4) \end{cases} \quad \text{\itshape (0,5 điểm)}$$
	\item Cộng từng vế của phương trình (3) và phương trình (4) ta được:
	$$(4x + 6y) + (9x - 6y) = 24 + 15 \iff 13x = 39 \iff x = 3. \quad \text{\itshape (0,75 điểm)}$$
	\item Thay $x = 3$ vào phương trình (1):
	$$2 \cdot 3 + 3y = 12 \iff 6 + 3y = 12 \iff 3y = 6 \iff y = 2. \quad \text{\itshape (0,5 điểm)}$$
	\item Kết luận: Vậy hệ phương trình đã cho có nghiệm duy nhất là $(x; y) = (3; 2)$. \quad \text{\itshape (0,25 điểm)}
\end{itemize}

\vspace{2mm}
\noindent\textbf{\large Bài 2 (3,0 điểm). Rút gọn biểu thức chứa căn thức:}

Cho biểu thức $Q = \left( \dfrac{\sqrt{x}}{\sqrt{x}-2} - \dfrac{4}{x-2\sqrt{x}} \right) : \left( \dfrac{1}{\sqrt{x}+2} + \dfrac{4}{x-4} \right)$ với $x > 0$ và $x \ne 4$.
\begin{itemize}
	\item Xét biểu thức trong ngoặc thứ nhất:
	$$A = \dfrac{\sqrt{x}}{\sqrt{x}-2} - \dfrac{4}{\sqrt{x}(\sqrt{x}-2)} = \dfrac{x - 4}{\sqrt{x}(\sqrt{x}-2)} = \dfrac{(\sqrt{x}-2)(\sqrt{x}+2)}{\sqrt{x}(\sqrt{x}-2)} = \dfrac{\sqrt{x}+2}{\sqrt{x}}. \quad \text{\itshape (1,25 điểm)}$$
	\item Xét biểu thức trong ngoặc thứ hai:
	$$B = \dfrac{1}{\sqrt{x}+2} + \dfrac{4}{(\sqrt{x}-2)(\sqrt{x}+2)} = \dfrac{\sqrt{x}-2+4}{(\sqrt{x}-2)(\sqrt{x}+2)} = \dfrac{\sqrt{x}+2}{(\sqrt{x}-2)(\sqrt{x}+2)} = \dfrac{1}{\sqrt{x}-2}. \quad \text{\itshape (1,0 điểm)}$$
	\item Thực hiện phép chia:
	$$Q = A : B = \dfrac{\sqrt{x}+2}{\sqrt{x}} : \dfrac{1}{\sqrt{x}-2} = \dfrac{\sqrt{x}+2}{\sqrt{x}} \cdot (\sqrt{x}-2) = \dfrac{x-4}{\sqrt{x}}.$$
	Vậy $Q = \dfrac{x-4}{\sqrt{x}}$. \quad \text{\itshape (0,75 điểm)}
\end{itemize}

\vspace{2mm}
\noindent\textbf{\large Bài 3 (3,0 điểm). Giải bài toán bằng cách lập hệ phương trình:}
\begin{itemize}
	\item Gọi số chi tiết máy tổ I phải làm theo kế hoạch là $x$ (chi tiết máy).\\
	Gọi số chi tiết máy tổ II phải làm theo kế hoạch là $y$ (chi tiết máy).\\
	Điều kiện: $x, y \in \mathbb{N}^*$ và $x, y < 800$. \quad \text{\itshape (0,5 điểm)}
	\item Vì theo kế hoạch hai tổ làm $800$ chi tiết máy nên ta có:
	$$x + y = 800 \quad (1) \quad \text{\itshape (0,5 điểm)}$$
	\item Thực tế tổ I vượt $20\%$ làm thêm $0{,}2x$; tổ II vượt $15\%$ làm thêm $0{,}15y$.\\
	Cả hai tổ làm được $945$ chi tiết máy (vượt mức $945 - 800 = 145$ chi tiết máy):
	$$0{,}2x + 0{,}15y = 145 \iff 20x + 15y = 14500 \iff 4x + 3y = 2900 \quad (2) \quad \text{\itshape (0,75 điểm)}$$
	\item Giải hệ:
	$$\begin{cases} x + y = 800 \\ 4x + 3y = 2900 \end{cases} \iff \begin{cases} 3x + 3y = 2400 \\ 4x + 3y = 2900 \end{cases} \iff \begin{cases} x = 500 \\ y = 300 \end{cases} \text{ (thỏa mãn điều kiện).} \quad \text{\itshape (0,75 điểm)}$$
	\item Kết luận: Theo kế hoạch tổ I làm $500$ chi tiết máy, tổ II làm $300$ chi tiết máy. \quad \text{\itshape (0,5 điểm)}
\end{itemize}

\vspace{5mm}
\begin{center}
\rule{5cm}{0.4pt}\\
{\small\itshape --- HẾT HƯỚNG DẪN CHẤM CHI TIẾT CẢ HAI MÃ ĐỀ ---}
\end{center}

\end{document}
